MOTOR SPECIFICATIONS (GM2804)
link: https://iflight-rc.eu/en-gb/products/ipower-gm2804-gimbal-motor-with-encoder
10.5Ohms between each M1, M2, M3 (when functional).
10V
0.8A load current
14N (7 pole pairs).
9 Ohms = phase resistance
K_v = 169 RPM/V
K_e = 0.057 V/rad/s
For our application we set V_d = 0 for simplicity (a good assumption at slow speeds.)
We calulate Vq = I * phase_resistance + V_bemf
where V_bemf = shaft_velocity / (K_v * SQRT(3)) / RPM_TO_RADS.


Explanation of FOC control
An N pole pair, XPWM BLDC motor has a stator with X coils around the outside and solid magnet with 2N "pizza slices", or poles (N, S, ...).
Each pole has a pair, so in our case we have a magnet pizza with 14 slices, made up of 7 opposing poles.
At any one time we can control the current through each of the coils independently.
The current through a coil creates a magnetic force in the direction that coil faces (see this great vid https://youtu.be/PgM8zWJr8-U?si=xZ6-wqovyPyhzzHs).
So we have three force vectors acting on the magnet pizza at any one time.
The direction of the force on the magnet at any one time is dependent not only on the currents on each of the coils but also the angular 
position of the pizza.
More specifically, the angle which matters to the force is the angle of the magnetic's magnetic field.
Because the pizza essentially repeats its North, South, North, South pattern N times, magnetic north repeats itself N times in one full mechanical
rotation.
So, to repeat, the force on the rotor (which is a pizza magnet) at any one time is determined by the angle of magnetic north of the pizza in
relation to the direction of the magnetic force generate by the three coils, as well as the magnitude of the force vector (ie the current) 
of each coil.
So lets give those things each a name.

Ia, Ib, Ic are the coil currents.
E is the angle between the magnetic fields, known as the electrical angle.

If we know the electrical angle, and we wish to generate a desired force vector, we can use the current electrical angle, along with the 
inverse Park and Clarke transforms to convert the desired force vector into the desired coil currents.

If we know the current current (lol) on the coild via current sensing, we can use PID control to decide the desired voltage.

We use PWM to communicate the desired voltages to the drv8313.


PID
The force vector we wish to generate will be based on the the velocity / angle / torque we wish to achieve.
If for instance we wish to achieve a certain mechanical angle, knowing our current one, we can use PID control to set estimate the desired
force vector.
The PID controller takes the error over time as an input, and calculates the derivative of that error at any one time (for a correction damping effect)
and the integral of the error over time (for a bias correction effect). The output of the PID controller is then
P * error + I * integral of error + D * derivative of error
where P, I, D are constants: the proportional, integral, and derivative gains respectively.

